\chapter{Conclusion}

\section[Summary]{\textbf{Summary}}

\subsection{Problem Addressed}

The EdgeCare Response platform was developed to address the critical gap between raw fall detection and clinically useful emergency response. While existing wearable monitors notify caregivers of a suspected fall, they fail to verify patient physiological stability, estimate emergency severity, coordinate transport, or prepare the receiving hospital. This leaves elderly patients vulnerable during the critical transport window.

\subsection{Objectives Achieved}

The project successfully achieved all stated objectives:
\begin{itemize}
\item A hybrid Edge-AI fall-prediction engine running on an ESP32 edge node, combining TinyML (1D CNN) inference with physical thresholds.
\item A physiological verification layer that normalized post-event heart rate against a rolling baseline to reduce false alarms.
\item An Emergency Severity Score (ESS) model to categorize events (Mild, Moderate, Critical) based on impact, physiology, inactivity, and user response.
\item An automated ambulance coordination and en route tracking workflow using Google Maps Directions API.
\item Real-time hospital pre-notification and cloud dashboard integration, enabling caregivers, dispatchers, and medical staff to monitor patient vitals and location simultaneously.
\end{itemize}

\subsection{Results Obtained}

The project successfully implemented an event-driven emergency healthcare system. The wearable edge node utilizes an ESP32, MPU6050 accelerometer/gyroscope, and a pulse sensor to sample data and perform local inference. Verified events are securely transmitted to a Next.js web application and Python FastAPI backend, which handles event logging and notifies stakeholders. Redis 7 serves as the real-time Pub/Sub message bus for WebSocket-based live dashboard updates, while PostgreSQL 15 handles persistent patient profiles, events, and response logs.

\subsection{Overall Outcome}

The platform provides a unique value proposition by linking wearable sensing to operational emergency coordination. EdgeCare Response successfully demonstrates that a low-power, privacy-aware wearable node combined with a scalable, low-latency cloud dashboard can significantly reduce emergency coordination times and false alarm rates, providing a solid foundation for clinical pilot studies.

\section[Personal Reflection]{\textbf{Personal Reflection}}

\subsection{Student 1 Reflection -- M P Preetham}

% M P Preetham - Personal reflection
Working on EdgeCare Response has been a transformative experience. Building the frontend dashboard using Next.js and WebSockets taught me the importance of real-time state management and responsive UI layout. Translating complex sensor status, alert severities, and live tracking data into an intuitive dashboard for caregivers, dispatchers, and hospitals showed me how software engineering can directly impact healthcare accessibility and response times.

\subsection{Student 2 Reflection -- Gagan B C}

% Gagan B C - Personal reflection
Developing the core event processing pipeline and TinyML inference for EdgeCare Response pushed my technical boundaries. Handling high-frequency sensor streams, implementing physiological verification logic, and ensuring sub-second delivery via Redis Pub/Sub highlighted the importance of low-latency, event-driven architectures. This project deepened my understanding of how local edge intelligence can work in harmony with cloud microservices.

\subsection{Student 3 Reflection -- Amaresh Malgi}

% Amaresh Malgi - Personal reflection
Designing the backend services and coordination workflows was both challenging and rewarding. Implementing the ambulance dispatch logic, PostgreSQL 15 database schema, and securing patient profile data gave me insights into secure, scalable distributed system design. EdgeCare Response allowed me to apply theoretical concepts of network communication and database transactions to a practical, life-saving application.

\subsection{Student 4 Reflection -- Gokul P}

% Gokul P - Personal reflection
Working on the hardware integration of the EdgeCare Response wearable node gave me hands-on experience with embedded systems and sensor interfacing. Connecting the MPU6050 and heart-rate sensor to the ESP32, calibrating inertial thresholds, and validating BLE data transmission taught me the critical importance of reliable hardware-software co-design in safety-critical applications. This project has solidified my understanding of IoT systems and their role in real-world healthcare solutions.
